\documentclass{article}
\usepackage{graphicx} % Required for inserting images
\usepackage{amsmath}
\usepackage{amssymb}
\usepackage{amsthm}
\usepackage{kotex}

\title{윌슨 정리에 기반한 새로운\\소수 판별 함수의 구성}
\author{이우진}
\date{}

\begin{document}

\maketitle

\section{서론}
소수 판별 함수란 1보다 큰 자연수 \(n\)에 대하여
\[\mathbf{1}_{\mathbb{P}}(n) = 
\begin{cases}
    1 & \text{(n은 소수)} \\
    0 & \text{(n은 합성수)}
\end{cases}\]
를 만족하는 함수 \(\mathbf{1}_{\mathbb{P}}\)를 말한다.
이 탐구 보고서에서는 윌슨 정리를 이용하여 새로운 소수 판별 함수를 구성한다.

\section{예비이론}
\subsection{나누어떨어짐}
정수 \(a \ne 0, b\)에 대하여 \(b\)가 \(a\)로 나누어떨어지면
\[a \mid b\]
로 표기한다.

\subsection{합동식}
정수 \(a, b\)와 자연수 \(m\)에 대하여 \(m \mid a-b\)일 때
\[a \equiv b \pmod m\]
이라고 표기하고, \(a\)와 \(b\)가 법 \(m\)에 대하여 합동이라고 한다.

합동식은 나눗셈 빼고 대부분의 등식의 성질을 만족한다.

\subsection{잉여역수}
정수 \(a\)와 자연수 \(m\)에 대하여 \(a\)와 \(m\)이 서로소이면 \[aa^{*} \equiv 1 \pmod m\] 인 정수 \(a^{*}\)가 존재하고, \(a^{*}\)를 \(a\)의 잉여역수라고 한다.

\subsection{최대정수함수}
실수 \(x\)에 대하여 정수 \(n\), 실수 \(\alpha\)가 \[x=n+\alpha, \quad 0 \le \alpha < 1\]을 만족할 때 \(\lfloor x \rfloor = n\)이다.

다른 말로 하면, \(\lfloor x \rfloor\)는 \(x\)의 내림이다.

\subsection{팩토리얼}
음이 아닌 정수 \(n\)에 대하여
\[n!=\prod_{k=1}^{n}k=1\times2\times3\times \cdots \times (n-1) \times n.\]
특히, \(0!\)은 빈 곱이므로 \(0!=1\)이다.

\section{윌슨 정리와 관련 보조정리}
\subsection{윌슨 정리}
1보다 큰 자연수 \(n\)에 대하여 \(n\)이 소수이면
\[(n-1)! \equiv -1 \pmod n\]
이고 그 역도 성립한다.

\begin{proof}
\((\rightarrow)\)

\(p\)가 소수라고 하자.
만약 \(p=2\)이면 \((2-1)!\equiv1\equiv-1\pmod 2\)이므로 성립한다.
그러므로 \(p\)가 홀수 소수라고 가정하자.
\(p\)보다 작은 모든 자연수는 \(p\)와 서로소이므로 각각 잉여역수를 가진다.
그 중 잉여역수가 자신인, 즉 \(aa \equiv a^2 \equiv1 \pmod p\)인 자연수 \(a\)는
\[p \mid a^2-1 = (a+1)(a-1)\]를 만족해야 한다.
\(p\)는 소수이므로 \(p \mid a + 1\) 또는 \(p \mid a - 1\)이다.
즉 \[a \equiv \pm 1 \pmod p\] 이다.

그러므로 \(1<k<p-1\)인 자연수 \(k\)는 자신과 합동이 아닌 잉여역수를 갖는다.
즉 \((p-1)! = 1 \times 2 \times \cdots \times (p-2) \times (p-1)\)에서 
\(2\)부터 \(p-2\)까지 곱하는 수의 개수는 \(p-3\)개로 짝수이므로 각각 잉여역수로 짝지어져 \(1\)과
합동이 되어 \[2 \times 3 \times \cdots (p-3) \times (p-2) \equiv 1 \pmod p\]가 된다.
그러므로 \[(p-1)! \equiv 1 \times (p-1) \times 2 \times 3 \times \cdots \times (p-2) \equiv 1 \times (p-1) \equiv -1 \pmod p\]이다.\\

\((\leftarrow)\)
정수 \(n\)이 \((n-1)! \equiv -1 \pmod n\)를 만족하는데 합성수라고 가정하자.
\(n\)은 합성수이므로 \(a \mid n, 1 < a < n\)인 자연수 \(a\)가 존재한다.
그러면 \(a\)는 다음 두 성질을 만족한다.
\begin{itemize}
    \item \(a \mid n \mid (n-1)! + 1\)
    \item \(a \le n-1\)이므로 \(a \mid (n-1)!\)
\end{itemize}

자연수 \(x, y, z\)에 대하여 \(x \mid y, x \mid z\)이면 \(x \mid y-z\)이다.
그러므로 \[a \mid (n-1)! + 1 - (n-1)! = 1.\]즉 \(a\)는 \(1\)의 양의 약수이므로 \(a = 1\)이어야 하는데 \(1 < a\)라고 가정했으므로 모순이 발생한다.
그러므로 \(n\)은 소수이다.
\end{proof}

\subsection{보조정리}
\(4\)보다 큰 합성수 \(n\)에 대하여
\[n \mid (n-1)!\]
이다.

\begin{proof}
먼저 \(n = ab, 1<a<b<n\)인 자연수 a, b가 존재하지 않는 경우를 생각하자. \(1\)이나 \(n\)이 아닌 \(n\)의 양의 약수가 하나만 존재해야 하므로 \(n=p^2\) (\(p\)는 소수)라고 할 수 있다. \(n>4\) 이므로 \(n > 2p\)이다.
따라서 \[n = p^2 \mid (n-1)! = 1 \times \cdots \times p \times \cdots \times 2p \times \cdots \times (n-1).\]

\(n = ab, 1<a<b<n\)인 자연수 a, b가 존재하는 경우에는
\[n = ab\mid (n-1)! = 1 \times \cdots \times a \times \cdots \times b \times \cdots \times (n-1).\]
\end{proof}

\section{새로운 소수 판별 함수의 구성}
일단 \(n \ne 4\)인 자연수만 생각하자. \(4\)가 아닌 합성수는 \(4\)보다 크므로 윌슨 정리와 보조정리 \(3.2\)에 의하여 
\[(n-1)! \equiv \begin{cases}
    n-1 & (n\text{은 소수})\\
    0 & (n\text{은 합성수})
\end{cases} \pmod n\]
이다. 여기에서 \((n-1)!\)을 \(n\)으로 나눈 나머지를 구하면 합동식의 결과를 얻을 수 있다. 

\(a, b, q, r\)이 모두 정수일 때 \(a \div b\)의 몫이 \(q\), 나머지가 \(r\)이면 \(a=bq+r\)이 성립한다. 
이항하면 \(a-bq=r\)이다. \(\left\lfloor \frac{a}{b} \right\rfloor\)는 \(a\)를 \(b\)로 나눈 몫이다. 

그러므로 구하는 나머지는
\[(n-1)! - n\left\lfloor \frac{(n-1)!}{n} \right\rfloor\]
이다. 여기에서 \(n\)이 소수일 때는 \(n-1\)이 아닌 \(1\)이 함수값이여야 하므로 전체 식을 \(n-1\)로 나누면 
\[(n-2)!-\frac{n}{n-1} \left\lfloor \frac{(n-1)!}{n} \right\rfloor\]
이다. 이제 이 함수는 \(4\)보다 큰 \(n\)에 대한 소수 판별 함수이다. 

\(n=4\)일 때는 보조정리가 성립하지 않기 때문에 추가 보정이 필요하다. 실제로 이 함수에 \(n=4\)를 대입하면
\[(4-2)!-\frac{4}{4-1} \left\lfloor \frac{(4-1)!}{4} \right\rfloor=2!-\frac{4}{3}=\frac{2}{3}\]
이라서 원하는 값인 \(0\)이 나오지 않는다. 그러므로 \(n=4\)일 때만 \(\frac{2}{3}\)을 빼줘야 한다. 
\(\left\lfloor \frac{1}{(n-4)^2+1} \right\rfloor\)는 \(n=4\)일 때만 \(1\)이고 다른 경우에는 \(0\)이다. 
이를 이용하여 최종적인 소수 판별 함수를 구하면
\[\mathbf{1}_\mathbb{P} (n)= (n-2)!-\frac{n}{n-1} \left\lfloor \frac{(n-1)!}{n} \right\rfloor\ - \frac{2}{3}\left\lfloor \frac{1}{(n-4)^2+1} \right\rfloor\]
이다.

\end{document}
